\documentclass[aspectratio=169, 14pt]{beamer}
\usepackage{stampfly_slides}

\lesson{11}{ホバリングタイム競技会}{Hover Time Competition}

\begin{document}

% ------------------------------------------------------------------
\begin{frame}
  \titlepage
\end{frame}

% ------------------------------------------------------------------
\begin{frame}{競技会概要 / Competition Overview}
  \begin{block}{ホバリングタイム競技}
    PID 制御で安定ホバリングし、滞空時間を競う
  \end{block}

  \vspace{1em}

  \begin{itemize}
    \item \textbf{種目}: ホバリングタイム（角速度フィードバック制御）
    \item \textbf{目標}: 最長 60 秒の安定ホバリング
    \item \textbf{評価}: 滞空時間（ベストタイム採用）
    \item \textbf{ツール}: \api{sf competition hover-time}
  \end{itemize}
\end{frame}

% ------------------------------------------------------------------
\begin{frame}{競技ルール詳細 / Competition Rules}
  \centering
  \begin{tabular}{ll}
    \hline
    \textbf{項目} & \textbf{内容} \\
    \hline
    種目  & ホバリングタイム（角速度フィードバック制御） \\
    制限時間 & 最長 60 秒 \\
    試行回数 & 3 回（ベストタイム採用） \\
    判定  & 離陸〜接地までの滞空時間 \\
    失格条件 & 安全エリア外への飛行、手による介入 \\
    \hline
  \end{tabular}

  \vspace{1.5em}

  \begin{alertblock}{注意}
    安全第一！\ 周囲に人がいないことを確認してから飛行
  \end{alertblock}
\end{frame}

% ------------------------------------------------------------------
\begin{frame}{タイムスケジュール / Schedule}
  \centering
  \begin{tabular}{ll}
    \hline
    \textbf{時間} & \textbf{内容} \\
    \hline
    \multicolumn{2}{l}{\textbf{Day 4（本日）}} \\
    13:00--13:30 & ルール説明（このスライド） \\
    13:30--16:00 & 準備・自由練習・予選 \\
    \hline
    \multicolumn{2}{l}{\textbf{Day 5（最終日）}} \\
    \phantom{0}9:00--\phantom{0}9:15 & ルール確認・機体チェック \\
    \phantom{0}9:15--\phantom{0}9:45 & 最終チューニング \\
    \phantom{0}9:45--10:45 & 競技本番 \\
    10:45--11:30 & 結果発表・表彰・振り返り \\
    \hline
  \end{tabular}
\end{frame}

% ------------------------------------------------------------------
\begin{frame}{攻略のヒント / Tips for Success}
  \begin{columns}[t]
    \begin{column}{0.48\textwidth}
      \begin{exampleblock}{PID ゲイン調整}
        \begin{itemize}
          \item まず P ゲインで安定化
          \item 振動には D ゲインを追加
          \item 定常偏差には I ゲイン調整
        \end{itemize}
      \end{exampleblock}
    \end{column}
    \begin{column}{0.48\textwidth}
      \begin{exampleblock}{バッテリ\,\&\,安全}
        \begin{itemize}
          \item フル充電で本番に臨む
          \item プロペラの取付を再確認
          \item 異常振動はすぐに着陸
        \end{itemize}
      \end{exampleblock}
    \end{column}
  \end{columns}

  \smallskip
  {\small \textbf{テレメトリ活用:}
    \api{sf log wifi} でリアルタイムデータを確認し、ゲインを追い込む}
\end{frame}

% ------------------------------------------------------------------
\begin{frame}{チェックポイント / Checkpoint}
  \begin{alertblock}{準備確認}
    \begin{itemize}
      \item[$\square$] 競技ルールと失格条件を理解した
      \item[$\square$] バッテリがフル充電されている
      \item[$\square$] PID ゲインの調整方法を把握している
      \item[$\square$] 安全チェックリストを確認した
    \end{itemize}
  \end{alertblock}

  \vspace{1em}

  \begin{block}{さあ、練習を始めよう！}
    13:30 から自由練習 → 最高のホバリングを目指そう
  \end{block}
\end{frame}

\end{document}
